% from ega4-2-fr.tex
\oldpage[IV-2]{86}

connexes de $X_{(\overline k)}$ ($1\leq i\leq m$); si
$p:X_{(K)}\to X_{(\overline k)}$ est la projection canonique, on sait que les
$p^{-1}(W_i)$ sont les composantes connexes de $X_{(K)}$ (4.4.6), donc
$\overline U$ est une réunion d'un certain nombre de ces composantes, et est
par suite défini sur $\overline k$; la conclusion résulte par suite de
(4.9.5, b)).
\end{proof}
\section{Dimension, profondeur, régularité dans les préschémas localement
noethériens}

Ce paragraphe se borne à reprendre, dans le langage géométrique, et avec
divers compléments de nature technique, les notions et résultats d'Algèbre
commutative exposés dans le chapitre 0, \S\S~16 et 17.

\subsection{Dimension des préschémas}

\begin{env}[5.1.1]
\label{IV.2.5.1.1-fr}
Soient $A$ un anneau, $\mathfrak J$ un idéal de $A$. Rappelons
$(0,16.1)$ que les définitions de la théorie de la dimension d'un anneau
$(0,16.1.1)$ et celles de la théorie de la dimension combinatoire des espaces
topologiques $(0,14.1.2)$ donnent les relations suivantes:
\begin{align}
\dim(\operatorname{Spec}(A))&=\dim(A),
\tag{5.1.1.1}\label{IV.2.5.1.1.1-fr}\\
\dim(V(\mathfrak J))&=\dim(A/\mathfrak J),
\tag{5.1.1.2}\label{IV.2.5.1.1.2-fr}\\
\operatorname{codim}(V(\mathfrak J),\operatorname{Spec}(A))
&=\operatorname{ht}(\mathfrak J).
\tag{5.1.1.3}\label{IV.2.5.1.1.3-fr}
\end{align}
\par\noindent\textup{(5.1.1.4)} $\operatorname{Spec}(A)$ est un espace
caténaire $\Longleftrightarrow A$ est un anneau caténaire.
\par\noindent\textup{(5.1.1.5)} $\operatorname{Spec}(A)$ est
équidimensionnel $\Longleftrightarrow A$ est équidimensionnel
$\Longleftrightarrow$ les idéaux premiers minimaux $\mathfrak p_\alpha$ de
$A$ sont tels que les dimensions des anneaux $A/\mathfrak p_\alpha$ soient
toutes égales.
\par\noindent\textup{(5.1.1.6)} $\operatorname{Spec}(A)$ est
équicodimensionnel $\Longleftrightarrow A$ est équicodimensionnel
$\Longleftrightarrow$ tous les idéaux maximaux de $A$ ont même hauteur.
On rappelle qu'un anneau noethérien $A$ est dit \emph{biéquidimensionnel} si
$\operatorname{Spec}(A)$ est biéquidimensionnel, c'est-à-dire si
$\operatorname{Spec}(A)$ est noethérien et $A$ est équidimensionnel,
équicodimensionnel, caténaire et de dimension finie.
\end{env}

\begin{proposition}[5.1.2]
\label{IV.2.5.1.2-fr}
Soient $X$ un préschéma, $Y$ une partie fermée irréductible de $X$, $y$ le
point générique de $Y$. On a
\begin{equation}
\operatorname{codim}(Y,X)=\dim(\mathcal O_{X,y}).
\tag{5.1.2.1}\label{IV.2.5.1.2.1-fr}
\end{equation}
\end{proposition}

\begin{proof}
En effet (I, 2.4.2) les parties fermées irréductibles de $X$ contenant $y$
sont canoniquement en correspondance biunivoque avec les parties fermées
irréductibles de $\operatorname{Spec}(\mathcal O_{X,y})$, donc sont en
correspondance biunivoque avec les idéaux premiers de $\mathcal O_{X,y}$, et
(5.1.2.1) résulte des définitions.

En particulier, on retrouve ainsi le fait que les points génériques des
composantes irréductibles de $X$ sont les \emph{seuls} points $x\in X$ tels
que $\dim(\mathcal O_x)=0$ (I, 1.1.14).
\end{proof}

